\pdfinfoomitdate=1
\pdftrailerid{}
\pdfsuppressptexinfo=15
\newcommand{\DoNotLoadEpstopdf}{}
\documentclass[11pt]{article}
\usepackage[T1]{fontenc}
\usepackage[utf8]{inputenc}
\usepackage{lmodern,amsmath,amssymb,amsthm,mathtools,booktabs,array}
\usepackage[margin=1in]{geometry}
\usepackage{microtype}
\usepackage[hidelinks]{hyperref}
\usepackage{enumitem}
\setlist{itemsep=3pt,topsep=5pt}
\allowdisplaybreaks[2]
\emergencystretch=2em
\newtheorem{theorem}{Theorem}[section]
\newtheorem{lemma}[theorem]{Lemma}
\newtheorem{proposition}[theorem]{Proposition}
\newtheorem{corollary}[theorem]{Corollary}
\theoremstyle{definition}\newtheorem{definition}[theorem]{Definition}
\theoremstyle{remark}\newtheorem{remark}[theorem]{Remark}
\newcommand{\R}{\mathbb R}
\newcommand{\Z}{\mathbb Z}
\newcommand{\Q}{\mathbb Q}
\newcommand{\T}{\mathbb T}
\newcommand{\E}{\mathbb E}
\newcommand{\Prob}{\mathbb P}
\newcommand{\covol}{\operatorname{covol}}
\newcommand{\tr}{\operatorname{tr}}
\newcommand{\Var}{\operatorname{Var}}
\newcommand{\Cov}{\operatorname{Cov}}
\newcommand{\dist}{\operatorname{dist}}
\newcommand{\ceilodd}{\operatorname{ceil}_{\rm odd}}
\newcommand{\normT}[1]{\lVert #1\rVert_{\T}}
\hypersetup{pdftitle={Report194 Asymptotics for Binary Matrices with Zero Row and Column Slopes},pdfauthor={Research report},pdfsubject={OEIS A222959 weighted binary matrices Fourier localization rational asymptotics}}
\title{Report194\\[5pt]Asymptotics for Binary Matrices\\with Zero Row and Column Slopes}
\author{A self contained study of the square diagonal of OEIS A222959}
\date{4 October 2026}
\begin{document}
\maketitle
\begin{abstract}
Let $C_n$ count the $n\times n$ binary matrices whose least-squares regression
slope is zero in every row and every column. With
$S_n=n(n^2-1)/12$ for odd $n$ and $S_n=n(n^2-1)/3$ for even $n$, we prove
\[
 C_n=2^{n^2}e^{-7/5}\sqrt{\frac2\pi}
       \left(\frac{2}{\pi S_n}\right)^{n-1}
       \left(1+\frac{171}{350n}+\frac{483051}{245000n^2}+O(n^{-3})\right).
\]
More generally, the factor in parentheses has an asymptotic expansion to every
fixed order, with rational coefficients independent of parity, given by a finite
Wick pairing and power-sum prescription. The main analytic issue is global
localization in a growing-dimensional Fourier integral. We prove a relative
exponential bound outside one embedded central tube by a dense arithmetic
progression inverse lemma and a defect-core decomposition. The proof includes
lattice saturation, Haar normalization, uniform core Gaussian ratios, and
uniform Gaussian polynomial moment estimates. We derive parity-aware smooth
inverse expansions and rigorous odd-integer threshold envelopes. Exact finite
counts and reproducible rational checks are supplied separately from the
asymptotic proof. The known small counts are strongly preasymptotic. No practical
onset, convergence of the full formal series, or absolute priority is claimed.
\end{abstract}

\paragraph{Attribution and scope.}
The matrix sequence is OEIS A222959 \cite{oeis959}. Its one-dimensional row
condition belongs to the earlier theory of fair binary words, including
Prodinger's asymptotics and Richomme's regression interpretation
\cite{prodinger,richomme}. The Fourier and higher-order Gaussian enumeration
strategy is part of the maximum-entropy and Edgeworth framework of Barvinok and
Hartigan \cite{bhgaussian,bhedgeworth}. The present report supplies a complete
problem-specific proof for the simultaneous growing arithmetic-progression row
and column constraints. The primary-source comparison in
Section~\ref{sec:literature} is deliberately scoped; a failure to locate an
identical theorem is not a certificate of originality.

\newpage
{\footnotesize\tableofcontents}
\newpage

\section{The counting problem and the main results}\label{sec:main}
All logarithms are natural. All asymptotic assertions refer to $n\to\infty$;
constants may depend on a fixed truncation order, but not on parity. The symbol
$O_J$ makes that dependence explicit. No constant below is optimized.

For $n\ge2$, regress a row $(x_1,\ldots,x_n)$ against the positions
$1,\ldots,n$, including an intercept. Its slope is
\[
 \frac{\sum_{j=1}^n(j-(n+1)/2)x_j}
      {\sum_{j=1}^n(j-(n+1)/2)^2}.
\]
Thus zero slope is a homogeneous weighted-sum condition. Define the primitive
integer weight vector $w=(w_1,\ldots,w_n)^T$ by
\begin{equation}\label{eq:weights}
 w_j=\begin{cases}j-(n+1)/2,&n\text{ odd},\\
                  2j-n-1,&n\text{ even}.
       \end{cases}
\end{equation}
The entries sum to zero, contain $1$, and form a symmetric arithmetic
progression of spacing $d_0=1$ or $2$. Write
\begin{equation}\label{eq:S}
 S=S_n=\sum_jw_j^2=
 \begin{cases}n(n^2-1)/12,&n\text{ odd},\\
 n(n^2-1)/3,&n\text{ even}.
 \end{cases}
\end{equation}
Then the square diagonal of A222959 is
\[
 C_n=\#\{M\in\{0,1\}^{n\times n}:Mw=M^Tw=0\}.
\]
We use $C_1=2$, the sequence's convention that its length-one rows and columns
are admissible; a regression slope denominator at $n=1$ is zero. Since
$\sum_jw_j=0$, replacing $M$ by the sign matrix $2M-\mathbf1\mathbf1^T$
preserves the two nullvector conditions. Consequently $C_n/2^{n^2}$ is the
probability that a random sign matrix has this prescribed vector as both a
right and a left nullvector.

Define
\begin{equation}\label{eq:G}
 G_n=\sqrt{\frac2\pi}\left(\frac2{\pi S_n}\right)^{n-1}
     =\left(\frac2\pi\right)^{n-1/2}S_n^{1-n}.
\end{equation}
The parity dependence in $S_n$ is essential: removing it would introduce an
exponential error.

\begin{theorem}[All fixed orders]\label{thm:main}
There is a sequence of rational numbers $(c_j)_{j\ge0}$, independent of parity,
with $c_0=1$, $c_1=171/350$, and $c_2=483051/245000$, such that, for every fixed integer $J\ge0$,
\begin{equation}\label{eq:main}
 C_n=2^{n^2}G_ne^{-7/5}
 \left(\sum_{j=0}^J\frac{c_j}{n^j}+O_J(n^{-J-1})\right).
\end{equation}
Section~\ref{sec:allorders} gives a finite exact coefficient prescription. No
convergence of the infinite formal series is asserted. The coefficients $c_0$, $c_1$, and $c_2$
are explicitly evaluated in this report.
\end{theorem}

\begin{theorem}[Relative global localization]\label{thm:globalintro}
Let $\T=\R/\Z$, let coefficient variables $a,b$ range over $\T^n$, and put
\[
 x_{ij}=w_ja_i+w_ib_j\pmod1,\qquad
 F(a,b)=\prod_{i,j}|\cos(\pi x_{ij})|.
\]
Let $E=\{aw^T+wb^T:a,b\in\R^n\}$ and
$\Lambda=E\cap\Z^{n\times n}$. Let $U$ be the image in $E/\Lambda$ of
$\{X\in E:\max_{ij}|X_{ij}|<1/16\}$. There is an absolute $c>0$ such that
\begin{equation}\label{eq:globalintro}
 \int_{\{(a,b):x\notin U\}}F(a,b)\,da\,db=O(G_ne^{-cn}).
\end{equation}
The measure on every coefficient torus is normalized Haar measure.
\end{theorem}

The second theorem is stronger than locating the exact maximum set of the
characteristic function. It controls all approximate modular resonances,
relative to the small main mass $G_n=\exp(-3n\log n+O(n))$.

\begin{corollary}[Parity and the first threshold]\label{cor:thresholdintro}
One has
\begin{equation}\label{eq:parityratio}
 \frac{C_{2m}}{C_{2m-1}}\sim\frac{6e^{-3}}{\pi m^3}.
\end{equation}
Each parity subsequence is eventually strictly increasing, whereas the
combined sequence is not. For all sufficiently large real $Y$, the first
threshold index
\[
 N(Y)=\min\{n\ge1:C_n\ge Y\}
\]
is odd. Section~\ref{sec:inverse} gives rigorous upper and lower envelopes for
$N(Y)$, including possible near-integer ambiguities.
\end{corollary}

\paragraph{Proof structure.}
Sections~\ref{sec:lattice} and \ref{sec:fourier} establish the exact lattice,
measure, and sign conventions. Sections~\ref{sec:inverselemmas}--\ref{sec:defects}
prove relative global localization. Sections~\ref{sec:gaussian}--\ref{sec:allorders}
then prove the fixed-order expansion, and Sections~\ref{sec:first} and~\ref{sec:second} compute its
first two corrections. The analytic estimates are independent of the finite checks
in Section~\ref{sec:checks}.

\section{Lattices and Gaussian normalization}\label{sec:lattice}
We work with Euclidean inner products, using the Frobenius inner product on
matrix spaces. The covolume of a full lattice in a subspace means the Euclidean
volume of one fundamental parallelepiped in that subspace.

\subsection{The margin lattice}
Let $A(M)=(Mw,M^Tw)$. Its integer image is exactly
\begin{equation}\label{eq:marginlattice}
 L=\{(r,c)\in\Z^{2n}:w^Tr=w^Tc\}.
\end{equation}
Indeed, the compatibility equation is necessary. Choose $u\in\Z^n$ with
$w^Tu=1$; one may choose a coordinate vector at a weight $1$. For any
compatible integer $(r,c)$, the integer matrix
\begin{equation}\label{eq:preimage}
 M=ru^T+uc^T-(w^Tr)uu^T
\end{equation}
satisfies $Mw=r$ and $M^Tw=c$. This proves saturation.

The primitive normal to the real margin subspace is $(w,-w)$, of length
$\sqrt{2S}$, so $\covol(L)=\sqrt{2S}$. To see the normal-vector formula,
complete the primitive normal to a unimodular integer basis, or equivalently
use a vector having integer dot product $1$ with it. The corresponding height
above its orthogonal hyperplane is the reciprocal of its length; unit ambient
volume is that height times the hyperplane lattice covolume.

For independent Bernoulli entries of mean $1/2$, the mean margin vector is zero
and its covariance on $H=(w,-w)^\perp$ is the restriction of
\[
 Q=\frac14\begin{pmatrix}S I_n&ww^T\\ww^T&S I_n\end{pmatrix}.
\]
Its nonzero eigenvalues are $S/4$, of multiplicity $2n-2$, and $S/2$, of
multiplicity one. Thus
\[
 \det{}'Q=\frac{2S^{2n-1}}{4^{2n-1}},\qquad
 \frac{\covol(L)}{(2\pi)^{(2n-1)/2}\sqrt{\det{}'Q}}=G_n.
\]
This is an exact calculation of the Gaussian lattice factor, not by itself a
local limit theorem.

\subsection{A rectangular phase lattice lemma}\label{sec:rectlattice}
The localization proof needs the same calculation after some rows and columns
are removed. The retained weights need not be centered.

\begin{lemma}[Primitive rectangular phase lattice]\label{lem:rectlattice}
Let $u\in\Z^r$ and $v\in\Z^s$ be nonzero primitive vectors, and write
$S_u=\|u\|^2$, $S_v=\|v\|^2$. Set
\[
 B(a,b)=av^T+ub^T,\quad E_{u,v}=\operatorname{im}_{\R} B,\quad
 \Lambda_{u,v}=E_{u,v}\cap\Z^{r\times s}.
\]
Then $\dim E_{u,v}=r+s-1$,
\begin{equation}\label{eq:rectcovol}
 B(\Z^{r+s})=\Lambda_{u,v},\qquad
 \covol(\Lambda_{u,v})=S_v^{(r-1)/2}S_u^{(s-1)/2}.
\end{equation}
The map from $\T^{r+s}$ to $E_{u,v}/\Lambda_{u,v}$ is onto and pushes normalized
Haar measure to normalized Haar measure. Its kernel is a single circle, with
no additional finite components.
\end{lemma}
\begin{proof}
The real kernel is $\R(u,-v)$, which proves the dimension claim. Choose integer
$p,q$ with $u^Tp=v^Tq=1$. For an integer $X\in E_{u,v}$, put
\[
 a=Xq,\qquad b=X^Tp-(p^TXq)v.
\]
If $X=a_0v^T+ub_0^T$, substitution gives $av^T+ub^T=X$; $a,b$ are integer.
This proves the first equality in \eqref{eq:rectcovol}.

The nonzero eigenvalues of $B^*B$ are
$S_v$ with multiplicity $r-1$, $S_u$ with multiplicity $s-1$, and
$S_u+S_v$ with multiplicity one. The integer kernel is generated by the
primitive vector $(u,-v)$, of length $K=\sqrt{S_u+S_v}$. Projection of
$\Z^{r+s}$ onto its orthogonal complement has covolume $1/K$: completing the
kernel vector to a unimodular basis gives unit volume equal to $K$ times the
projected base volume. The Jacobian of $B$ on that complement is
\[
 \sqrt{S_v^{r-1}S_u^{s-1}(S_u+S_v)}.
\]
Multiplying by $1/K$ gives the claimed covolume.

Surjectivity on the real quotient is immediate, and the integer image
calculation makes the torus map well-defined. An onto homomorphism of compact
groups pushes normalized Haar measure to normalized Haar measure: its
pushforward is a probability measure invariant under every target translation.
Finally, if $B(a,b)$ is an integer matrix, saturation gives an integer pair
$(a',b')$ with the same image. Hence $(a-a',b-b')$ lies in the real kernel.
Modulo integer coefficients, the kernel is exactly that real line modulo its
primitive integer generator, a circle.
\end{proof}

On any embedded real subset of $E_{u,v}/\Lambda_{u,v}$ the normalized Haar
density is $1/\covol(\Lambda_{u,v})$ times Euclidean volume. In particular,
\begin{equation}\label{eq:rectG}
 \frac1{\covol(\Lambda_{u,v})}
 \int_{E_{u,v}}e^{-\pi^2\|X\|^2/2}\,dX
 =\left(\frac2\pi\right)^{(r+s-1)/2}
   S_v^{-(r-1)/2}S_u^{-(s-1)/2}.
\end{equation}
For $u=v=w$, this is exactly $G_n$. The quotient calculation explains why no
extra gauge-circle length belongs in the answer.

\section{Fourier inversion and the positive central tube}\label{sec:fourier}
The elementary identity $\int_{\T}e^{2\pi ikt}\,dt=1_{\{k=0\}}$ yields
\begin{equation}\label{eq:fourier}
 \frac{C_n}{2^{n^2}}=
 \int_{\T^{2n}}\prod_{i,j}
    \frac{1+e^{2\pi i(w_ja_i+w_ib_j)}}2\,da\,db.
\end{equation}
One may integrate all $2n$ margins even though one equation is redundant.
For real representatives of the coefficient variables, write
$X_{ij}=w_ja_i+w_ib_j$. Because the full weight lists are centered,
$\sum_{ij}X_{ij}=0$. Thus the signed characteristic function is
\begin{equation}\label{eq:Phi}
 \Phi(a,b)=\prod_{ij}\cos(\pi X_{ij}),\qquad |\Phi|=F.
\end{equation}
Although one cosine is not a function on $\R/\Z$, its absolute value is, and
the full signed product is well-defined. An integer change of a coefficient
changes its sign by $(-1)^{\sum_jw_j}=1$.

For $0<r<1/2$, let $U_r$ consist of the phase-torus classes admitting a
representative $X\in E$ with $\max|X_{ij}|<r$. Such a representative is unique
when $2r<1$: two representatives differ by an integer matrix with every entry
of absolute value less than one. We use $r=1/16$ and write $U=U_{1/16}$.
If an arbitrary full real representative is $X+K$, with $X$ in this tube and
$K\in\Lambda$, then $\sum K_{ij}=0$, since every matrix in $E$ has total sum
zero. Consequently
\[
 \prod\cos(\pi(X_{ij}+K_{ij}))
 =(-1)^{\sum K_{ij}}\prod\cos(\pi X_{ij})
 =\prod\cos(\pi X_{ij})>0.
\]
There is therefore no local sign cancellation or missing parity multiplier.

The only exact maximum phases of $F$ are integer phases, which constitute the
identity of $E/\Lambda$; Lemma~\ref{lem:rectlattice} identifies their coefficient
preimage as the gauge circle $(a,b)=(tw,-tw)$. This fact alone does not exclude
nearby or distant approximate resonances. The next three sections establish
the necessary quantitative statement.

We will use the global scalar estimate
\begin{equation}\label{eq:cosbound}
 |\cos(\pi t)|\le
 \exp\left(-\frac{\pi^2}2\normT{t}^2\right),
 \qquad \normT{t}=\dist(t,\Z).
\end{equation}
For $|x|<1/2$, the derivative of $\log\cos(\pi x)+\pi^2x^2/2$ is
$-\pi\tan(\pi x)+\pi^2x$, nonpositive for $x\ge0$; symmetry and the endpoint
limit prove the inequality. Outside the central tube, the modular distance in
\eqref{eq:cosbound} cannot be replaced by an arbitrary real lift.

\section{Dense arithmetic progression inverse lemmas}\label{sec:inverselemmas}
All statements in this section hold for sufficiently large $n$, uniformly in
its parity. Their deliberately generous constants make the required density
and no-wrap conditions transparent.

\begin{lemma}[Scalar inverse]\label{lem:scalar}
Suppose $0<\tau<1/8$ and $0\le\rho<1/4$. If $t\in\T$ satisfies
$\normT{w_jt}\le\tau$ for at least $(1-\rho)n$ indices, then
\[
 \normT{t}\le\frac{8\tau}{n}.
\]
\end{lemma}
\begin{proof}
Let $H=\lfloor n/2\rfloor$. For each $1\le h\le H$, at least one pair
$j,j+h$ has both endpoints good: there are $n-h$ pairs and at most $2\rho n$
are spoiled. Subtraction gives $\normT{d_0ht}\le2\tau$ for every such $h$.

We first note an elementary no-wrap fact. If $\normT{hz}\le A<1/4$ for
$1\le h\le H$, then $\normT{z}\le A/H$. Otherwise, since
$\normT{z}\le A$, the first $h$ with $h\normT{z}>A$ satisfies
$A<h\normT{z}\le2A<1/2$. At this first crossing there is no wrap, contradicting
$\normT{hz}\le A$.

Apply this fact with $z=d_0t$ and $A=2\tau$. If $d_0=1$, it gives the result.
If $d_0=2$, then $t$ is within $\tau/H$ of an integer or a half-integer. The
half-integer branch is impossible: every weight is odd, and any good weight
would have distance at least $1/2-n\tau/H\ge1/2-3\tau>\tau$ from an integer.
The integer branch gives the claimed bound. The constants cover all rounding
for sufficiently large $n$.
\end{proof}

\begin{lemma}[Dense circuit stability]\label{lem:dense}
Put $\kappa=10^{-4}$. Let $I\subseteq\{1,\ldots,n\}$ omit at most
$\kappa n$ indices. Suppose $\theta_i\in\T$, $i\in I$, and
\begin{equation}\label{eq:circuit}
 \normT{\sum_i z_i\theta_i}\le\eta
\end{equation}
for every integer vector $z$ supported in $I$ with
$\sum_i z_iw_i=0$ and $\sum_i|z_i|\le4$. If $6n\eta<1/2$, there is $t\in\T$
such that
\begin{equation}\label{eq:denseconclusion}
 \normT{\theta_i-w_it}\le500\eta\qquad(i\in I).
\end{equation}
\end{lemma}
\begin{proof}
Set $H=\lfloor n/4\rfloor$. For each $0\le h\le H$, choose a pair
$x,x+h\in I$ and define $f(h)=\theta_{x+h}-\theta_x$, with $f(0)=0$.
Pairs exist because $n-h-2\kappa n>0$. Any two choices differ in circle norm by
at most $\eta$: their difference is a four-term circuit in \eqref{eq:circuit}.
For $h,k\ge0$, $h+k\le H$, a triple $x,x+h,x+h+k\in I$ exists because at
least $n-H-3\kappa n>0$ starting positions remain. Comparing the three chosen
increments gives
\begin{equation}\label{eq:fadd}
 \normT{f(h+k)-f(h)-f(k)}\le3\eta.
\end{equation}

Choose $F(0)=0$, choose any real lift $F(1)$, and choose subsequent real lifts
recursively with $|F(h)-F(h-1)-F(1)|\le3\eta$. Then
$|F(h)-hF(1)|\le3(h-1)\eta$. For positive $h,k$ with $h+k\le H$, the real
additive defect has absolute value at most $6H\eta<1/2$. It is a lift of the
circle element in \eqref{eq:fadd}, so its absolute value is at most $3\eta$.
This step is what prevents an accumulated integer winding error.

Put $A=F(H)/H$ and $g(h)=F(h)-Ah$. For $2h\le H$,
$|2g(h)-g(2h)|\le3\eta$. For $2h>H$, add the additive relations for
$H=h+(H-h)$ and $h=(H-h)+(2h-H)$ to obtain
$|2g(h)-g(2h-H)|\le6\eta$. Since $g(0)=g(H)=0$, the maximum
$M=\max_{0\le h\le H}|g(h)|$ satisfies $M\le M/2+3\eta$, hence
$M\le6\eta$. Every actual pair in $I$ at distance at most $H$ consequently
satisfies
\[
 \normT{\theta_y-\theta_x-A(y-x)}\le7\eta.
\]
Partition the full index interval into sixteen consecutive blocks of nearly
equal length. Every block contains an element of $I$. For sufficiently large
$n$, adjacent selected representatives are at distance at most $H$, and any
point can be joined to one fixed representative by at most sixteen such links.
Thus, for some real $b$,
\[
 \normT{\theta_i-Ai-b}\le112\eta\qquad(i\in I).
\]
Recenter the progression: in circle notation
$\theta_i=(A/d_0)w_i+B+e_i$, with real lifts $|e_i|\le112\eta$.

If $n=2m+1$ is odd, view the retained weights as a dense subset of
$[-m,m]\cap\Z$. Choose a retained $p$ in
$[-\lfloor m/4\rfloor,\lfloor m/4\rfloor]$. Among the possible $q$ in
$[-\lfloor m/2\rfloor,\lfloor m/2\rfloor]$, at most $2\kappa n$ fail one of
the conditions that $q$ and $p+q$ be retained. All sums are inside $[-m,m]$,
so a retained triple $p,q,r=p+q$ exists. Its three-term circuit gives
$\normT{B}\le\eta+3(112\eta)=337\eta$. Repeated indices simply combine
coefficients and cause no problem. Hence the total error is at most
$449\eta<500\eta$.

If $n$ is even, a retained opposite-weight pair exists, since otherwise at
least $n/2$ weights would be omitted. It gives
$\normT{2B}\le\eta+2(112\eta)=225\eta$. Thus $B$ is within $112.5\eta$ of
an integer or a half-integer. All $w_i$ are odd, so a half-integer constant is
absorbed into the coefficient of $w_i$ modulo one. The total error is then at
most $224.5\eta<500\eta$. This proves \eqref{eq:denseconclusion} in both
parities.
\end{proof}

\section{A robust central core and its Gaussian cost}\label{sec:core}
Fix, once and for all,
\begin{equation}\label{eq:smallconstants}
 \epsilon=10^{-3},\qquad\kappa=10^{-4},\qquad\delta=10^{-10}.
\end{equation}
An entry is called bad when $\normT{x_{ij}}\ge\delta$.

\begin{lemma}[Robust central core]\label{lem:core}
Let $R,C\subseteq\{1,\ldots,n\}$ each omit at most $\kappa n$ indices.
Suppose every row in $R$ and every column in $C$ has at most $\epsilon n$ bad
entries in its full row or column. There is $t\in\T$ such that
\begin{align}
 \normT{a_i-w_it}&\le16000\delta/n&& (i\in R),\label{eq:corea}\\
 \normT{b_j+w_jt}&\le2064064\delta/n&& (j\in C).\label{eq:coreb}
\end{align}
The core phases have a real lift in $E_{w_R,w_C}$ with every entry of absolute
value at most $2080064\delta<1/16$.
\end{lemma}
\begin{proof}
For any integer circuit $z$ supported in $R$ with $\ell^1$ norm at most four,
at most $4\epsilon n$ full columns meet a bad entry in its support. In every
other column,
\[
 \normT{w_j\sum_i z_i a_i}
 =\normT{\sum_i z_i x_{ij}}\le4\delta.
\]
Lemma~\ref{lem:scalar}, with $\rho=4\epsilon<1/4$ and $\tau=4\delta$, gives
$\normT{\sum_i z_i a_i}\le32\delta/n$. Apply Lemma~\ref{lem:dense} with
$\eta=32\delta/n$; its lift condition follows from $192\delta<1/2$. We get
$a_i=w_it+e_i$ modulo one, where $|e_i|\le16000\delta/n$. By the transposed
argument $b_j=w_js+f_j$, with $|f_j|\le16000\delta/n$.

Put $h=t+s$. On a good core entry, using $|w_i|\le n$,
\[
 \normT{w_iw_jh}\le\delta+|w_je_i|+|w_if_j|
 \le32001\delta=:\tau.
\]
For each retained row more than $3n/4$ retained columns are good. Scalar
inversion in $j$ gives $\normT{w_ih}\le8\tau/n$ for every $i\in R$.
The retained row set is dense, so a second scalar inversion gives
\[
 \normT{h}\le64\tau/n^2=2048064\delta/n^2.
\]
Choose the nearest-integer real lift of $h$ and combine it with $f_j$. This
gives \eqref{eq:coreb}, while \eqref{eq:corea} already holds. Using these real
residual coefficients, the lift is $e_iw_j+w_i\widetilde f_j$; it genuinely
lies in the core phase space and its maximum is at most
$(16000+2064064)\delta=0.0002080064<1/16$.
\end{proof}

For retained sizes $r=|R|$, $s=|C|$, write
\[
 S_R=\sum_{i\in R}w_i^2,\quad S_C=\sum_{j\in C}w_j^2,\quad
 G_{R,C}=\left(\frac2\pi\right)^{(r+s-1)/2}
 S_C^{-(r-1)/2}S_R^{-(s-1)/2}.
\]
A dense retained list contains consecutive indices. For odd $n$, the
corresponding weights are consecutive integers; for even $n$, they are
consecutive odd numbers with difference two. Their greatest common divisor is
one in either case. Thus both retained lists are primitive and
Lemma~\ref{lem:rectlattice} applies.

On the embedded core tube, \eqref{eq:cosbound} and \eqref{eq:rectG} give
\begin{equation}\label{eq:coreintegral}
 \int_{\{\text{core in its central tube}\}}
 \prod_{i\in R,j\in C}|\cos(\pi x_{ij})|\,da_R\,db_C\le G_{R,C}.
\end{equation}
Only an absolute-product estimate is used here. Restricted weights need not
sum to zero, so we do not assert the signed-integrand property for a core.

\begin{lemma}[Uniform Gaussian cost per defect]\label{lem:ratio}
Let $k=n-r$, $l=n-s$, and $v=k+l\le\kappa n$. Then
\begin{equation}\label{eq:coreratio}
 \frac{G_{R,C}}{G_n}\le(C_Gn^{3/2})^v,
 \qquad C_G=\sqrt{\pi/2}\,e^{16}.
\end{equation}
\end{lemma}
\begin{proof}
Uniformly in both parities, $n^3/16\le S\le n^3$ for $n\ge2$, and
$\max_iw_i^2\le n^2$. Therefore
\[
 S_R/S\ge1-16k/n,\quad S_C/S\ge1-16l/n,
\]
and, as $16\kappa<1/2$,
\[
 \log(S/S_R)\le32k/n,\qquad\log(S/S_C)\le32l/n.
\]
The exact ratio is
\[
 \frac{G_{R,C}}{G_n}=
 \left(\frac\pi2\right)^{v/2}S^{v/2}
 \left(\frac S{S_C}\right)^{(r-1)/2}
 \left(\frac S{S_R}\right)^{(s-1)/2}.
\]
Bounding the last two logarithms by $16l$ and $16k$ proves the claim.
In particular there is no extraneous factor $e^{O(n)}$ independent of $v$.
\end{proof}

\section{Global relative localization by defects}\label{sec:defects}
We now prove Theorem~\ref{thm:globalintro}. Define the modular energy
\[
 \mathcal E(a,b)=\sum_{i,j}\normT{x_{ij}}^2,
 \qquad c_0=\delta^2\epsilon\kappa/2.
\]
On $\mathcal E>c_0n^2$, \eqref{eq:cosbound} gives the integral bound
\begin{equation}\label{eq:highenergy}
 \int_{\{\mathcal E>c_0n^2\}}F\le e^{-\pi^2c_0n^2/2}
 =O(G_ne^{-cn})
\end{equation}
for some $c>0$ and all sufficiently large $n$, because
$\log G_n=-3n\log n+O(n)$. Comparison with $G_n$ is essential here.

On the low-energy region, remove simultaneously every row and every column
with more than $\epsilon n$ bad entries, counting bad entries in the original
full rows and columns. Let $k,l$ be the numbers removed, $v=k+l$, and let
$R,C$ be the retained sets. If $B$ is the total number of bad entries, then
$B\delta^2\le\mathcal E$, and
\begin{equation}\label{eq:defectnumber}
 v\le\frac{2B}{\epsilon n}
 \le\frac{2\mathcal E}{\delta^2\epsilon n}\le\kappa n.
\end{equation}
Thus the retained core satisfies Lemma~\ref{lem:core}.

Each removed row has at least $\epsilon n-l\ge\epsilon n/2$ bad entries in
retained columns. Each removed column likewise has at least $\epsilon n/2$
bad entries in retained rows, since $\kappa\le\epsilon/2$. These charges are
on the disjoint sets $R^c\times C$ and $R\times C^c$, and no charged factor is
counted twice. There are at least $\epsilon nv/2$ charged factors, so
\begin{equation}\label{eq:crosspenalty}
 F\le e^{-\pi^2\delta^2\epsilon nv/4}
       \prod_{i\in R,j\in C}|\cos(\pi x_{ij})|.
\end{equation}
This analytic form avoids evaluating $\cos(\pi\delta)$ in floating point,
where it would round to one.

Fix particular removed sets, and let $\mathcal A_{R,C}$ be the event that these
are exactly the removed sets and that $\mathcal E\le c_0n^2$. Put
$\gamma=\pi^2\delta^2\epsilon/4>0$. On this event the core lies in its central
tube. Enlarge just that event after applying \eqref{eq:crosspenalty}:
\begin{align}
 \int 1_{\mathcal A_{R,C}}F\,da\,db
 &\le e^{-\gamma nv}
 \int1_{\{\text{core in central tube}\}}
       \prod_{i\in R,j\in C}|\cos(\pi x_{ij})|\,da\,db\notag\\
 &\le e^{-\gamma nv}G_{R,C}.\label{eq:patternintegral}
\end{align}
The discarded coefficient coordinates integrate over full tori of total
volume one. The remaining variables push forward to core Haar measure as in
Lemma~\ref{lem:rectlattice}. This argument does not assume independence of the
defect event and does not introduce a conditional fiber-volume loss.

For each $v$, the number of pairs of removed sets is
$\sum_{k+l=v}\binom nk\binom nl=\binom{2n}v\le(2n)^v$. By
\eqref{eq:coreratio}, the total contribution of nonempty low-energy defect
patterns is at most
\begin{equation}\label{eq:defectsum}
 G_n\sum_{v\ge1}\left(2C_G n^{5/2}e^{-\gamma n}\right)^v
 =O(G_ne^{-cn}).
\end{equation}
The last step holds eventually, because the bracket is at most
$e^{-\gamma n/2}$ for all sufficiently large $n$. If $v=0$,
Lemma~\ref{lem:core} applies to the whole phase matrix and puts it in $U$.
Hence every point outside $U$ is covered by \eqref{eq:highenergy} or
\eqref{eq:defectsum}. This proves Theorem~\ref{thm:globalintro}.

\paragraph{What the constants do and do not establish.}
The small fixed $\delta$ proves existence of a positive localization rate and
an eventual threshold. It is not a usable numerical error estimate. The cost
of one defect is only polynomial in $n$, whereas its cross-edge penalty is
exponential in $n$; that comparison is the structural reason the proof closes.

\section{The local Gaussian and a dimension free moment bound}\label{sec:gaussian}
\subsection{Projection and shrinking cutoff}
Put $v_i=w_i/\sqrt S$, $a_i=v_i^2$, and $s_r=\sum_i a_i^r$. The orthogonal
projection onto $E$ is
\begin{equation}\label{eq:Paction}
 P(X)=Xvv^T+vv^TX-vv^TXvv^T.
\end{equation}
Indeed the two first summands are commuting orthogonal projections whose
intersection is spanned by $vv^T$. In entry notation,
\begin{equation}\label{eq:P}
 P_{(ij),(kl)}=1_{\{i=k\}}v_jv_l+1_{\{j=l\}}v_iv_k-v_iv_jv_kv_l.
\end{equation}
Its diagonal entries are
\begin{equation}\label{eq:h}
 h_{ij}=a_i+a_j-a_ia_j,\qquad \max h_{ij}=O(n^{-1}),\quad
 \sum_{ij}h_{ij}=2n-1.
\end{equation}
Here $s_1=1$ and $\max a_i=O(n^{-1})$, directly from \eqref{eq:S}.

Under the normalized Gaussian with density proportional to
$e^{-\pi^2\|X\|^2/2}$ on $E$, the matrix $Y=\pi X$ is a centered Gaussian with
covariance $P$. Fix any $0<\zeta<1/6$ and let
\[
 \mathcal B_n=\{\max_{ij}|Y_{ij}|\le n^{-1/2+\zeta}\}.
\]
The Gaussian one-coordinate tail and a union bound give
\begin{equation}\label{eq:gausstail}
 \Prob(\mathcal B_n^c)\le2n^2e^{-c_1n^{2\zeta}}.
\end{equation}
For completeness, the scalar bound follows by applying Markov's inequality to
$e^{tZ}$, using $\E e^{tZ}=e^{t^2\sigma^2/2}$, and minimizing over $t$; apply
it to both signs. The cutoff tube lies in $U$ for sufficiently large $n$.
On the intervening part of $U$, the absolute integrand is bounded by its
Gaussian quadratic part. Global localization, the positive sign in $U$, and
\eqref{eq:gausstail} therefore imply, for every fixed $M$,
\begin{equation}\label{eq:Gaussianreduction}
 \frac{C_n}{2^{n^2}G_n}=
 \E\left[1_{\mathcal B_n}
  \exp\left\{\sum_{ij}\left(\log\cos Y_{ij}+\frac{Y_{ij}^2}2\right)\right\}\right]
 +O_M(n^{-M}).
\end{equation}
All logs inside this cutoff are of positive cosines.

\subsection{A self contained Gaussian polynomial estimate}
We record the dimension-free estimate required to expand the quartic term.
For a random variable $Z$, write $\|Z\|_p=(\E|Z|^p)^{1/p}$.

\begin{lemma}[Fixed-degree Gaussian moments]\label{lem:moments}
For each fixed nonnegative integer $d$ and $1\le p<\infty$, there is a constant
$C_{d,p}$ independent of the dimension such that every polynomial $F(g)$ of
degree at most $d$ in a standard Gaussian vector satisfies
\[
 \|F\|_p\le C_{d,p}\|F\|_2.
\]
\end{lemma}
\begin{proof}
Choose an even integer $2m\ge p$; norm monotonicity reduces the claim to $2m$.
Define the multivariate Wick polynomials by the generating identity
\[
 e^{t^Tg-\|t\|^2/2}=\sum_{k\ge0}\frac1{k!}
 \sum_{i_1,\ldots,i_k}t_{i_1}\cdots t_{i_k}
 :g_{i_1}\cdots g_{i_k}:.
\]
Taking the expectation of products of these generating functions yields
$\exp(\sum_{a<b}t_a^Tt_b)$. Coefficient comparison proves orthogonality of
different total degrees and Wick's rule: pairings join slots in different
Wick-ordered factors only. Since these polynomials have the ordinary monomials
as leading terms, every $F$ has an orthogonal finite decomposition
$F=\sum_{k=0}^dF_k$, where
\[
 F_k=\sum_{i_1,\ldots,i_k}T_{i_1\ldots i_k}
 :g_{i_1}\cdots g_{i_k}:
\]
for a symmetric tensor $T$, and $\|F_k\|_2^2=k!\|T\|_{\rm HS}^2$.

For $k\ge1$, the moment $\E F_k^{2m}$ is a sum over at most $(2mk-1)!!$
pairings of slots, with no pairing internal to one factor. Each pairing is a
tensor contraction along a loop-free multigraph with $2m$ tensor vertices.
Its absolute value is at most $\|T\|_{\rm HS}^{2m}$. To prove this last claim,
if tensors $A,B$ share index tuple $I$ and have disjoint remaining tuples $J,K$,
form $C_{J,K}=\sum_I A_{I,J}B_{I,K}$. Cauchy--Schwarz gives
\[
 \sum_{J,K}|C_{J,K}|^2
 \le\left(\sum_{I,J}|A_{I,J}|^2\right)
      \left(\sum_{I,K}|B_{I,K}|^2\right).
\]
Merge two connected tensor blocks and contract all their common edges.
Continue in each connected component; all newly internal edges are removed at
each merge, so no uncontrolled self-trace appears. Disconnected components
multiply. This proves the contraction bound and hence
\[
 \|F_k\|_{2m}\le
 \left(\frac{(2mk-1)!!}{(k!)^m}\right)^{1/(2m)}\|F_k\|_2.
\]
The $k=0$ case is immediate. Minkowski followed by Cauchy--Schwarz over the
fixed set $0\le k\le d$ and orthogonality proves the claim, with no ambient
dimension in the constant.
\end{proof}

\subsection{Quartic concentration and higher powers}
Define positive rational coefficients by
\[
 -\log\cos z=\frac{z^2}2+\sum_{k\ge2}q_kz^{2k},\qquad
 q_2=\frac1{12},\quad q_3=\frac1{45},\quad q_4=\frac{17}{2520},
\]
and put $Q_{2k}=q_k\sum_{ij}Y_{ij}^{2k}$. Rationality and positivity follow
recursively from $\tan'=1+\tan^2$, $\tan0=0$, followed by integration.
For each fixed $k,p$, Minkowski and scalar Gaussian moments give
\begin{equation}\label{eq:Qmoments}
 \|Q_{2k}\|_p\le C_{k,p}\sum_e h_e^k
 \le C_{k,p}(\max h)^{k-1}\tr P=O_{k,p}(n^{2-k}),
\end{equation}
where a single index $e$ denotes a cell.

The quartic mean is
\begin{equation}\label{eq:m}
 m(n)=\E Q_4=\frac14\sum_eh_e^2
      =\frac{2ns_2+2-4s_2+s_2^2}{4},\qquad
 s_2=\frac{3(3n^2-7)}{5n(n^2-1)}.
\end{equation}
In particular $m(n)=7/5+O(n^{-1})$. For jointly centered Gaussian variables,
Wick pairing gives
\[
 \Cov(Y_e^4,Y_f^4)=72h_eh_fP_{ef}^2+24P_{ef}^4.
\]
There are $72$ pairings with exactly two cross pairs and $24$ with four;
the no-cross pairings cancel in the covariance. Thus
\begin{equation}\label{eq:variance}
 \Var Q_4=\frac12\sum_{e,f}h_eh_fP_{ef}^2+
          \frac16\sum_{e,f}P_{ef}^4=O(n^{-1}).
\end{equation}
Indeed $|P_{ef}|^2\le h_eh_f$, by positivity of a covariance matrix, and
$\sum_{e,f}P_{ef}^2=\tr P=2n-1$. Each displayed sum is bounded by a constant
times $(\max h)^2\tr P$.
For $D=Q_4-m(n)$, Lemma~\ref{lem:moments} consequently yields
\begin{equation}\label{eq:Dmoments}
 \|D\|_p=O_p(n^{-1/2})
\end{equation}
for every fixed $p$. This remains valid in the growing dimension $2n-1$.

\section{A constructive expansion to every fixed order}\label{sec:allorders}
We now prove Theorem~\ref{thm:main}, except for the evaluations of $c_1$ and $c_2$.

\subsection{Two finite Taylor expansions with controlled remainders}
Fix $K\ge2$. On $|y|\le1/2$, positivity and analyticity of the log-cosine
series imply
\[
 0\le-\log\cos y-y^2/2-\sum_{k=2}^Kq_ky^{2k}
 \le C_K|y|^{2K+2}.
\]
Hence on $\mathcal B_n$, for all sufficiently large $n$,
\[
 \sum_e(\log\cos Y_e+Y_e^2/2)=-\sum_{k=2}^KQ_{2k}-R_K,
 \qquad0\le R_K\le C_K\sum_e|Y_e|^{2K+2}.
\]
The difference of the two exponentials with and without $R_K$ is at most
$R_K$, since the retained $Q$'s are nonnegative. Its expected value is
$O_K(\sum_e h_e^{K+1})=O_K(n^{1-K})$. Removing the cutoff from
$\exp(-\sum Q_{2k})\le1$ costs at most \eqref{eq:gausstail}. Therefore
\begin{equation}\label{eq:finiteQ}
 \frac{C_n}{2^{n^2}G_n}=\E\exp\left(-\sum_{k=2}^KQ_{2k}\right)
 +O_K(n^{1-K}).
\end{equation}
No pointwise maximum bound multiplied by $n^2$ is used for this remainder;
the integrated Gaussian moments provide its stated order.

For the target order $J\ge0$, choose
\begin{equation}\label{eq:KLZ}
 K=J+3,\qquad L=2J+3,\qquad Z_K=D+\sum_{k=3}^KQ_{2k}.
\end{equation}
Equations~\eqref{eq:Qmoments} and \eqref{eq:Dmoments} show
$\|Z_K\|_p=O_{J,p}(n^{-1/2})$. Crucially,
$Z_K=\sum_{k=2}^KQ_{2k}-m(n)\ge-m(n)$. All derivatives of $e^{-z}$ along the
segment from $0$ to $Z_K$ are bounded by $e^{m(n)}=O(1)$. The ordinary finite
Taylor theorem, for every realization, thus gives
\[
 \left|e^{-Z_K}-\sum_{r=0}^L\frac{(-Z_K)^r}{r!}\right|
 \le\frac{e^{m(n)}}{(L+1)!}|Z_K|^{L+1}.
\]
Its expectation is $O_J(n^{-J-2})$, as is the error in \eqref{eq:finiteQ}.
Define the finite quantity
\begin{equation}\label{eq:TJ}
 T_J(n)=\sum_{r=0}^{2J+3}\frac{(-1)^r}{r!}\E Z_{J+3}^r.
\end{equation}
We have proved the stronger intermediate approximation
\begin{equation}\label{eq:finiteapproximant}
 \frac{C_n}{2^{n^2}G_n}=e^{-m(n)}T_J(n)+O_J(n^{-J-2}).
\end{equation}
There has been no exchange of an infinite series with an expectation. The
one-sided bound on $Z_K$ avoids any need for an exponential moment of $|D|$.

The polynomial families in this argument also have negligible weighted tails
on $\mathcal B_n^c$: H\"older and the fixed-moment bounds give polynomial
$L^2$ bounds for every fixed product, and Cauchy--Schwarz with
\eqref{eq:gausstail} makes the tail superpolynomially small. This assertion is
about these controlled families, not arbitrary polynomials with unrestricted
$n$-dependent coefficients.

\subsection{Exact rationality by Wick pairings and power sums}
For this part use $u_i=i-(n+1)/2$ for every $n$, even when the $u_i$ are
half-integers, and put $S_u=n(n^2-1)/12$. Scaling all weights by a nonzero
constant does not change $E$ or $P$. Thus in both parities the same covariance
formula is
\begin{equation}\label{eq:rationalP}
 P_{(ij),(kl)}=
 1_{\{i=k\}}\frac{u_ju_l}{S_u}
 +1_{\{j=l\}}\frac{u_iu_k}{S_u}
 -\frac{u_iu_ju_ku_l}{S_u^2}.
\end{equation}
Every fixed Wick moment in \eqref{eq:TJ} is an exact rational function of $n$.
Here is a constructive reduction proving that claim.
\begin{enumerate}
\item Expand the finitely many powers of $D=Q_4-m(n)$ and the higher $Q$'s.
All entry-index sums initially run without restrictions; $m(n)$ is already
rational.
\item Apply Wick's rule to each Gaussian monomial, a finite sum of covariance
products. The number of slots and pairings depends only on $J$.
\item For each covariance factor choose one of the three terms in
\eqref{eq:rationalP}. Merge index variables wherever a Kronecker delta imposes
equality.
\item Sum each remaining free index independently. Every result is a rational
coefficient times a power of $S_u^{-1}$ and a product of centered power sums
\begin{equation}\label{eq:Mr}
 M_r(n)=\sum_{i=1}^n\left(i-\frac{n+1}2\right)^r,\qquad r\ge0.
\end{equation}
\end{enumerate}
Each $M_r(n)$ is a polynomial with rational coefficients. One elementary
power-sum recursion is
\[
 \sum_{i=1}^n i^r=
 \frac{(n+1)^{r+1}-1-\sum_{k=0}^{r-1}\binom{r+1}k\sum_{i=1}^n i^k}{r+1},
\]
which follows by telescoping $(i+1)^{r+1}-i^{r+1}$. Binomially expand the
centered powers in \eqref{eq:Mr}; odd centered powers vanish by reflection,
and $M_0(n)=n$. This supplies a finite rational implementation without invoking
an unevaluated summation theorem.

Repeated entry indices cause no missing case. They are present in the initial
unrestricted sums, and the delta equalities only merge variables. No condition
that surviving indices be distinct occurs, so no hidden inclusion-exclusion
is required. This proves exact rationality and parity independence of $T_J$.
Primitive integral weights remain essential for the lattice factor $G_n$;
they are not required for the covariance calculation.

\subsection{Coefficient extraction and consistency}
Because $\E D=0$,
$\E Z_K=\sum_{k=3}^K\E Q_{2k}=O(n^{-1})$. For every fixed $r\ge2$,
$|\E Z_K^r|\le\|Z_K\|_r^r=O(n^{-r/2})$. Hence
\[
 T_J(n)=1+O(n^{-1}).
\]
As a rational function, $T_J$ therefore has no pole at infinity: its Laurent
expansion contains only a constant and negative powers of $n$. Likewise
$m(n)-7/5=O(n^{-1})$ has rational coefficients there. For $0\le j\le J$,
define
\begin{equation}\label{eq:cj}
 c_j=[z^j]\left\{\exp\bigl(7/5-m(1/z)\bigr)T_J(1/z)\right\}.
\end{equation}
The expression in braces is analytic near $z=0$, has constant term one, and
has rational Taylor coefficients. Its finite Taylor remainder through degree
$J$ is $O_J(n^{-J-1})$ at $z=1/n$. Together with
\eqref{eq:finiteapproximant}, this proves \eqref{eq:main}.

If two target orders $J,J'$ are used, both finite analytic approximants
approximate the same normalized exact count. Their difference along $z=1/n$
is $O(n^{-\min(J,J')-2})$. An analytic function with a first nonzero term of
lower degree could not have that bound. Their coefficients therefore agree
through every common displayed order. Thus \eqref{eq:cj} defines one consistent
sequence $(c_j)$.

This is a constructive all-fixed-orders result, not an efficient algorithm
claim: naive Wick enumeration is expensive. It does not imply that the full
formal series converges, that $J$ may grow with $n$, or that truncation errors
are numerically small at the currently tabulated dimensions.

\section{The explicit first correction}\label{sec:first}
We give enough exact algebra to isolate the value $171/350$ and to make its
independent rational checks transparent. Finite centered power summation gives
\begin{align}
 s_2&=\frac{3(3n^2-7)}{5n(n^2-1)}=\frac9{5n}+O(n^{-3}),\label{eq:s2}\\
 s_3&=\frac{9(3n^4-18n^2+31)}{7n^2(n^2-1)^2}
       =\frac{27}{7n^2}+O(n^{-4}).\label{eq:s3}
\end{align}
These identities hold exactly for both parities because normalized weights are
scale-invariant. Equation~\eqref{eq:m} gives
\begin{equation}\label{eq:mexp}
 m(n)=\frac75-\frac9{5n}+O(n^{-2}).
\end{equation}
Expanding $(a_i+a_j-a_ia_j)^3$ and summing yields
\begin{align}
 H_3:=\sum_{ij}h_{ij}^3
 &=2ns_3+6s_2-6s_3-6s_2^2+6s_2s_3-s_3^2\notag\\
 &=\frac{648}{35n}+O(n^{-2}).\label{eq:H3}
\end{align}
Since $\E Y_e^6=15h_e^3$,
\begin{equation}\label{eq:EQ6}
 \E Q_6=\frac13H_3=\frac{216}{35n}+O(n^{-2}).
\end{equation}

\subsection{The two variance sums}
Write
\[
 A=\sum_{e,f}h_eh_fP_{ef}^2,\qquad B=\sum_{e,f}P_{ef}^4.
\]
Let $U,V$ be respectively the same-row and same-column projections in
\eqref{eq:P}, and let $R=zz^T$ with $z=v\otimes v$. Then $P=U+V-R$ and
$A=\tr(HPHP)$ for $H=\operatorname{diag}(h)$. The following intermediate
finite sums evaluate its trace terms:
\begin{align*}
 F_0&=\sum_i[a_i+(1-a_i)s_2]^2=3s_2+(n-4)s_2^2+s_2^3,\\
 T_0&=\sum_{ij}h_{ij}^2a_ia_j=2s_3+s_3^2+2s_2^2-4s_2s_3,\\
 W_0&=\sum_i a_i[a_i+(1-a_i)s_2]^2
       =(1-s_2)^2s_3+2s_2^2(1-s_2)+s_2^2,\\
 V_0&=\sum_{ij}h_{ij}a_ia_j=2s_2-s_2^2.
\end{align*}
Expansion of the trace gives $A=2F_0+2T_0-4W_0+V_0^2$. Hence the exact identity
\begin{equation}\label{eq:Aexact}
 A=6s_2+(2n-12)s_2^2+6s_2^3+s_2^4-4s_2^2s_3+2s_3^2
  =\frac{432}{25n}+O(n^{-2}).
\end{equation}
For $B$, divide the ordered pair of cells into four disjoint cases: same row
only, same column only, equal cells, and different rows and columns. With
$s_4=\sum_i a_i^4$, this gives
\begin{equation}\label{eq:Bpartition}
 B=2\left[\sum_i(1-a_i)^4\right](s_2^2-s_4)
     +\sum_{ij}h_{ij}^4+(s_2^2-s_4)^2.
\end{equation}
For example, on the same-row-only case,
$P_{(ij),(il)}=(1-a_i)v_jv_l$, and the sum over $j\ne l$ contributes
$s_2^2-s_4$. Direct expansion gives
\begin{align*}
 \sum_{ij}h_{ij}^4={}&6s_2^2-24s_2s_3+12s_2s_4+12s_3^2-8s_3s_4\\
 &+8s_3+s_4^2+2ns_4-8s_4.
\end{align*}
After substitution into \eqref{eq:Bpartition}, every $s_4$ term cancels:
\begin{align}
 B={}&(2n-2)s_2^2+12s_2^3+s_2^4+8s_3-24s_2s_3
             -8s_2^2s_3+12s_3^2\notag\\
   ={}&\frac{162}{25n}+O(n^{-2}).\label{eq:Bexact}
\end{align}
By \eqref{eq:variance},
\begin{equation}\label{eq:varlead}
 \Var Q_4=\frac12A+\frac16B=\frac{243}{25n}+O(n^{-2}).
\end{equation}

\subsection{Remainder control and the coefficient}
Taking $K=3$ in \eqref{eq:finiteQ} gives an $O(n^{-2})$ error. Since
$D\ge-m$ and $m=O(1)$, Taylor's theorem and \eqref{eq:Dmoments} yield
\[
 \E e^{-Q_4}=e^{-m}\left(1+\frac12\Var Q_4+O(n^{-3/2})\right).
\]
Moreover $\|Q_6\|_2=O(n^{-1})$ and $e^{-x}$ is $1$-Lipschitz on
$[0,\infty)$, so
\[
 \left|\E[(e^{-Q_4}-e^{-m})Q_6]\right|
 \le\|D\|_2\|Q_6\|_2=O(n^{-3/2}).
\]
Using $|e^{-x}-1+x|\le x^2/2$ for $x\ge0$, we get
\[
 \E e^{-Q_4-Q_6}
 =e^{-m}\left(1+\frac12\Var Q_4-\E Q_6+O(n^{-3/2})\right).
\]
The order $n^{-1}$ correction inside these parentheses is
\[
 \frac{243}{50}-\frac{216}{35}=-\frac{459}{350}.
\]
Multiplication by \eqref{eq:mexp} contributes $9/5$, giving
\[
 c_1=\frac95-\frac{459}{350}=\frac{171}{350}.
\]
This calculation on its own has an $O(n^{-3/2})$ remainder. The all-fixed-orders
theorem already proved in Section~\ref{sec:allorders}, and uniqueness of the
first coefficient, improve the first-order statement to $O(n^{-2})$. No
uncomputed higher coefficient is inferred from this remainder improvement.

\section{The explicit second correction}\label{sec:second}
The finite coefficient prescription can be evaluated one order further. We
prove
\begin{equation}\label{eq:c2}
 c_2=\frac{483051}{245000},\qquad
 \log\left(\frac{e^{7/5}C_n}{2^{n^2}G_n}\right)
 =\frac{171}{350n}+\frac{6483}{3500n^2}+O(n^{-3}).
\end{equation}
No coefficient is fitted to finite counts. The additional ingredient is a
connected-Wick power bound, followed by two finite contractions.

\subsection{Connected Wick pairings and their power count}
For random variables with moments, joint cumulants are defined recursively by
\[
 \E\prod_{i=1}^r Z_i=
 \sum_{\mathcal P}\prod_{B\in\mathcal P}\kappa(Z_i:i\in B),
\]
where the sum is over all set partitions of $\{1,\ldots,r\}$. This is a finite
algebraic definition and does not assume existence of a moment-generating
function near zero.

\begin{lemma}[Connected power count]\label{lem:cumulants}
For fixed integers $k_i\ge2$,
\begin{equation}\label{eq:cumulantbound}
 \kappa(Q_{2k_1},\ldots,Q_{2k_r})
 =O\left(n^{1-\sum_{i=1}^r(k_i-1)}\right).
\end{equation}
\end{lemma}
\begin{proof}
Represent each Gaussian monomial by a vertex with $2k_i$ slots. Wick pairings
give multigraphs on these vertices, allowing self-loops. Grouping pairings by
their connected components is exactly the moment-cumulant recursion, so the
cumulant is the sum of connected pairings.

There are $M=\sum_i k_i$ paired edges. In each covariance from \eqref{eq:P},
choose the row-equality term, the column-equality term, or the negative product
term, denoted respectively R, C, and X. Let $t$ be the number of X edges, and
let $\rho_R,\rho_C$ be the ranks of the R and C edge graphs on the $r$ vertices;
a self-loop has rank zero. Kronecker equalities leave
$b=2r-\rho_R-\rho_C$ free row and column index blocks. The total exponent of
all normalized weights $v_i$ is $2M+2t$. A free block of exponent $d\ge0$ has
absolute sum at most $O(n^{1-d/2})$, since $\max_i|v_i|=O(n^{-1/2})$;
this includes an unweighted block with $d=0$. The colored term therefore is
\[
 O(n^{2r-\rho_R-\rho_C-M-t}).
\]
Replace the R and C subgraphs by spanning forests in each of their components,
and retain the X edges. Connectivity of the full graph is preserved, whence
$\rho_R+\rho_C+t\ge r-1$. The exponent is at most $r+1-M$, which is the
one in \eqref{eq:cumulantbound}. Only finitely many pairings and color choices
occur at fixed degrees. This argument bounds their absolute sums, without
using a favorable cancellation or a sign assumption on covariance entries.
\end{proof}

\subsection{Two finite contractions}
Write $x=s_2$, $y=s_3$, $z=s_4$, and use
\[
 F=Q_6,\quad H=Q_8,\quad V=\E D^2,\quad
 K=\Cov(Q_4,Q_6),\quad T=\kappa(Q_4,Q_4,Q_4)=\E D^3.
\]
These local letters are unrelated to the phase characteristic function or the
margin subspace used earlier. The bivariate Wick formula gives
\begin{equation}\label{eq:Kcontraction}
 K=\sum_{e,f}\left(h_eh_f^2P_{ef}^2+\frac23h_fP_{ef}^4\right).
\end{equation}
For $T$, write
$D=\sum_e(\tfrac12h_e:Y_e^2:+\tfrac1{12}:Y_e^4:)$.
The four possible triples of chaos degrees give
\begin{align}
 T={}&\sum_{e,f,g}h_eh_fh_gP_{ef}P_{eg}P_{fg}
 +\frac32\sum_{e,f,g}h_eh_fP_{eg}^2P_{fg}^2\notag\\
 &+2\sum_{e,f,g}h_eP_{ef}P_{eg}P_{fg}^3
 +\sum_{e,f,g}P_{ef}^2P_{eg}^2P_{fg}^2.\label{eq:Tcontraction}
\end{align}
All cell indices are unrestricted. For example, the chaos degree triples
$(2,2,2)$, $(2,2,4)$, $(2,4,4)$, and $(4,4,4)$ have respectively
$8$, $24$, $192$, and $1728$ pairings for one ordering; the coefficients in
\eqref{eq:Tcontraction} include the three possible positions in the two mixed
cases and the coefficients in $D$.

For reproducibility, the complete polynomial reductions of these contractions
are as follows:
\begin{align}
 K={}&\frac{10}3nxy+\frac{34}3x^2+4y
 -\frac{52}3x^3-2x^4+\frac{16}3x^3y+\frac{148}3x^2y
 -\frac53x^2y^2\notag\\
 &-\frac{40}3x^2z-\frac{40}3xy-\frac{76}3xy^2
 +\frac{28}3xyz-\frac{64}3xz-22y^2+40yz+\frac{16}3z-10z^2,
 \label{eq:Kexact}\\
 T={}&11nx^3+37x^2+2y-104x^3+124x^4+40x^5+\frac{11}2x^6
 -50x^4y-232x^3y\notag\\
 &+70x^3z-114x^2y+94x^2y^2+50x^2z+46xy+258xy^2
 -210xyz\notag\\
 &+8xz-22y^2-102yz+4z+90z^2.\label{eq:Texact}
\end{align}
We specify an exact finite derivation so these are not identities inferred
from interpolation. For a degree vector $(d_1,\ldots,d_r)$, enumerate loop
counts $\ell_i$ and nonnegative off-diagonal edge multiplicities $e_{ij}$ with
$2\ell_i+\sum_{j\ne i}e_{ij}=d_i$, retaining connected graphs. The number of
slot pairings for each graph is
\begin{equation}\label{eq:graphmultiplicity}
 \frac{\prod_i d_i!}
      {\prod_i(2^{\ell_i}\ell_i!)\prod_{i<j}e_{ij}!}.
\end{equation}
This follows by assigning the slots at each vertex to its incident edges,
pairing its $2\ell_i$ loop slots, and matching slots across each pair of
vertices. For every edge multiplicity $e$, choose $a+b+c=e$ R, C, X colors,
with coefficient $(-1)^c e!/(a!b!c!)$. Merge row blocks along R edges and
column blocks along C edges. At each endpoint an R edge adds one column
weight power, a C edge one row power, and an X edge one of each; a loop adds
twice at its endpoint. After merging, a block of even degree $2j$ contributes
$s_j$ and a block of odd degree contributes zero, by reflection. Here $s_0=n$.
Multiply by \eqref{eq:graphmultiplicity} and the appropriate $q_k$'s, and sum.
This finite distributive expansion of \eqref{eq:P} proves
\eqref{eq:Kexact}--\eqref{eq:Texact} for every $n\ge2$. The package implements
exactly this symbolic computation and checks the two polynomial coefficient
lists independently.

As a shorter check of the terms actually needed asymptotically, there are two
connected graphs for $(4,6)$ and eight for $(4,4,4)$. Their contributions are
shown below; each last-column entry is per labeled graph and omits only
$O(n^{-3})$ terms, including all $q_k$ factors.
\begin{center}
\small
\begin{tabular}{p{2.35in}rrl}
\toprule
Degree pattern and edges&Graphs&Pairings&Contribution per graph\\
\midrule
$(4,6)$, two joining edges, loops $1,2$&1&540&$2nxy+4y+10x^2$\\
$(4,6)$, four joining edges, loops $0,1$&1&360&$\frac43nxy+\frac43x^2$\\
$(4,4,4)$, double-edge star, leaf loops&3&864&$nx^3+7x^2$\\
$(4,4,4)$, triangle, one loop each&1&1728&$2nx^3+2y+12x^2$\\
$(4,4,4)$, edges $1,1,3$, one loop&3&1152&$\frac43nx^3+\frac43x^2$\\
$(4,4,4)$, double triangle, no loops&1&1728&$2nx^3$\\
\bottomrule
\end{tabular}
\end{center}
It follows, also directly from the exact polynomials, that
\begin{align}
 K&=\frac{10}3nxy+4y+\frac{34}3x^2+O(n^{-3})
    =\frac{13176}{175n^2}+O(n^{-3}),\label{eq:Klead}\\
 T&=11nx^3+2y+37x^2+O(n^{-3})
    =\frac{167778}{875n^2}+O(n^{-3}).\label{eq:Tlead}
\end{align}

\subsection{A finite exponential expansion with an order three remainder}
Equation~\eqref{eq:finiteQ} with $K=4$ gives
\[
 \frac{C_n}{2^{n^2}G_n}=\E e^{-Q_4-F-H}+O(n^{-3}).
\]
Set $Z=D+F+H$. Its fixed moments obey
$\|Z\|_p=O_p(n^{-1/2})$ and $Z\ge-m$. Taylor expansion of $e^{-Z}$ through
degree five therefore has expected remainder $O(\E|Z|^6)=O(n^{-3})$.
To control every monomial in that finite polynomial more sharply than by its
absolute moments, use the cumulant partitions of $\E D^aF^bH^c$.
Singleton $D$ blocks vanish. Every surviving partition has at most
$\lfloor a/2\rfloor+b+c$ blocks: blocks containing no $F,H$ require at least
two $D$'s, and each other block can be charged to one $F$ or $H$. By
Lemma~\ref{lem:cumulants}, its product is
\[
 O\bigl(n^{-a-2b-3c+\#\text{blocks}}\bigr)
 =O\bigl(n^{-\lceil a/2\rceil-b-2c}\bigr).
\]
In the second equality the right side is an upper bound; a partition can have
fewer than the maximal number of blocks. Centering changes no cumulant of
order at least two. In particular,
\[
 \E D^4=3V^2+O(n^{-3}),\quad
 \E D^2F=V\E F+O(n^{-3}),\quad
 \E F^2=(\E F)^2+O(n^{-3}).
\]
Keeping the terms of orders at most two gives, with $\mu_F=\E F$ and
$\mu_H=\E H$,
\begin{align}
 \E e^{-Q_4-F-H}=e^{-m}\bigg[&1-\mu_F-\mu_H+\frac V2+K+
 \frac{\mu_F^2}2-\frac T6\notag\\
 &-\frac{V\mu_F}2+\frac{V^2}8+O(n^{-3})\bigg].\label{eq:secondexp}
\end{align}
All other degree-at-most-five terms are $O(n^{-3})$ by the displayed partition
bound. No infinite cumulant series has been interchanged with an expectation.

\subsection{Rational assembly}
Centered power summation gives
\[
 x=\frac9{5n}-\frac{12}{5n^3}+O(n^{-5}),\quad
 y=\frac{27}{7n^2}+O(n^{-4}),\quad z=\frac9{n^3}+O(n^{-5}).
\]
More generally, $s_r=n^{1-r}(3^r/(2r+1)+O(n^{-2}))$ for fixed $r$; this
follows by expanding centered even power sums, or by the midpoint sum for the
fixed polynomial of degree $2r$. The mean of $Q_8$ is
$\mu_H=(17/24)\sum_eh_e^4$, using the Gaussian eighth moment $105$.
Equations~\eqref{eq:m}, \eqref{eq:Aexact}, \eqref{eq:Bexact},
\eqref{eq:H3}, and the displayed fourth leverage sum give the coefficients
\begin{center}
\begin{tabular}{crr}
\toprule
Quantity&Coefficient of $n^{-1}$&Coefficient of $n^{-2}$\\
\midrule
$m-7/5$&$-9/5$&$-39/100$\\
$V$&$243/25$&$-2691/175$\\
$\mu_F$&$216/35$&$-2484/175$\\
$\mu_H$&$0$&$8466/175$\\
$K$&$0$&$13176/175$\\
$T$&$0$&$167778/875$\\
\bottomrule
\end{tabular}
\end{center}
Taking the ordinary logarithm of the finite expansion
\eqref{eq:secondexp} cancels its square terms. Its first coefficient is
$L_1=171/350$ as before, and the second is
\[
 L_2=\frac{39}{100}-\frac{2691}{350}+\frac{2484}{175}
       -\frac{8466}{175}+\frac{13176}{175}-\frac{167778}{5250}
     =\frac{6483}{3500}.
\]
Exponentiating this finite expansion gives
\[
 c_2=L_2+\frac{L_1^2}2
     =\frac{6483}{3500}+\frac{29241}{245000}
     =\frac{483051}{245000},
\]
proving \eqref{eq:c2}. The symbolic contraction computation and the independent
raw-Gaussian-moment checks in the package use rational arithmetic; they do not
use finite matrix counts to estimate this coefficient.

\section{Parity aware inversion and integer thresholds}\label{sec:inverse}
\subsection{The logarithmic phase}
Use $d=12$ on odd indices and $d=3$ on even indices, so that
$S_n=n(n^2-1)/d$. Define
\begin{equation}\label{eq:inverseconstants}
 A_0=\log2,\quad B_0=\log(2/\pi),\quad
 D_d=\log(2d/\pi),\quad E_d=-\log d-B_0/2-7/5.
\end{equation}
Taking logarithms in Theorem~\ref{thm:main} gives, on the specified parity,
\begin{equation}\label{eq:logphase}
 \log C_n=A_0n^2-3n\log n+D_dn+3\log n+E_d
          +\sum_{j=1}^Jr_jn^{-j}+O_J(n^{-J-1}).
\end{equation}
The $r_j$ are rational and independent of parity. They are the formal
coefficients of
\[
 -(n-1)\log(1-n^{-2})+
       \log\left(1+\sum_{j\ge1}c_jn^{-j}\right),
\]
so in particular $r_1=1+c_1=521/350$ and $r_2=6483/3500-1=2983/3500$.
Each requested finite order involves only
finitely many coefficients. Positivity needed for the logarithm follows from
the leading equivalent.

\subsection{Why the first crossing is eventually odd}
The leading equivalent, evaluated at adjacent parities, gives
\[
 \frac{C_{2m}}{C_{2m-1}}\sim
 2^{4m-1}\frac2\pi\,\frac1{S_{2m}}
       \left(\frac{S_{2m-1}}{S_{2m}}\right)^{2m-2}.
\]
Here
\[
 S_{2m}=\frac{2m(2m-1)(2m+1)}3,\qquad
 \frac{S_{2m-1}}{S_{2m}}=\frac{m-1}{2(2m+1)}.
\]
Canceling the powers of four leaves \eqref{eq:parityratio}, since
$((2m-2)/(2m+1))^{2m-2}\to e^{-3}$.
On either fixed parity,
\begin{equation}\label{eq:step2ratio}
 \frac{C_{n+2}}{C_n}\sim
 2^{4n+4}e^{-6}\left(\frac{2d}{\pi}\right)^2n^{-6}\longrightarrow\infty.
\end{equation}
Thus each parity subsequence eventually increases strictly and tends to
infinity, while every sufficiently large even term is below its preceding odd
term. Once $Y$ exceeds the maximum of the finitely many exceptional early
terms, its first crossing cannot be even. This proves
Corollary~\ref{cor:thresholdintro}, without asserting a numerical onset.

\subsection{The inverse of a smooth asymptotic phase}
Let $L=\log Y$ and $X=\sqrt{L/A_0}$. A large real root of a smooth truncation of
\eqref{eq:logphase}, for either $d$, has the expansion
\begin{align}
 x={}&X+u+
 \frac{A_0u^2+3u-3\log X-E_d}{2A_0X}
 +O\left(\frac{\log^3X}{X^2}\right),\label{eq:smoothinverse}\\
 u={}&\frac{3\log X-D_d}{2A_0}.\notag
\end{align}
To verify it, substitute $x=X+u+v/X$. The order-$X$ terms cancel when
$2A_0u=3\log X-D_d$. The remaining constant terms are
$2A_0v-A_0u^2-3u+3\log X+E_d$, which gives the displayed $v$. All omitted
phase terms are $O(\log^3X/X)$; the phase derivative is asymptotic to
$2A_0X$, giving the stated inverse error. Higher inverse coefficients follow
recursively by formal substitution, using the finite-order remainders to
justify any fixed truncation. The $r_1/x$ term first affects a later inverse
order than the $X^{-1}$ term displayed here.

\subsection{Rigorous odd integer envelopes}
For the odd branch $d=12$, let $M_J(x)$ be the right side of
\eqref{eq:logphase} truncated after $r_Jx^{-J}$, omitting the sum if $J=0$.
There are constants $K_J,n_J$ such that, for all odd $n\ge n_J$,
\begin{equation}\label{eq:phasebounds}
 M_J(n)-K_Jn^{-J-1}\le\log C_n\le M_J(n)+K_Jn^{-J-1}.
\end{equation}
Both bounding functions are strictly increasing for all sufficiently large
real $x$, since their derivatives are asymptotic to $2A_0x$. Let $x_+(L)$ and
$x_-(L)$ be their large roots:
\[
 M_J(x_+)+K_Jx_+^{-J-1}=L,\qquad
 M_J(x_-)-K_Jx_-^{-J-1}=L.
\]
Then $x_+\le x_-$. Define
\[
 \ceilodd(x)=2\left\lceil\frac{x-1}2\right\rceil+1,
\]
the least odd integer not below $x$. For all sufficiently large $Y$,
\begin{equation}\label{eq:oddenvelope}
 \ceilodd(x_+(\log Y))\le N(Y)\le\ceilodd(x_-(\log Y)).
\end{equation}
Indeed every large odd $n<x_+$ has its upper phase bound below $L$, while the
odd integer $\ceilodd(x_-)$ has its lower phase bound at least $L$. The finite
initial segment is excluded by increasing $Y$ if necessary. This proof also
covers an exact root at an odd integer, because the bounds use $\ge$ at a
threshold.

The mean value theorem gives
\[
 x_-(L)-x_+(L)=O_J(X^{-J-2}).
\]
Away from an odd integer within this shrinking interval, the two ceilings
agree and determine $N(Y)$ exactly. When they straddle an odd integer, the
asymptotic information intentionally leaves a rounding ambiguity. No uniform
single rounded smooth formula follows at all thresholds.

\paragraph{An interpolation qualification.}
Formula~\eqref{eq:smoothinverse} concerns a smooth asymptotic phase. Even on
one parity branch, ordinary piecewise-linear interpolation across intervals
of length two has a chord error of order one, since the smooth phase's second
derivative tends to $2\log2$. This can produce an inverse error of order
$X^{-1}$, precisely the scale of the displayed second correction. Interpolation
between adjacent indices of different parities has the additional branch
oscillation. Neither interpolation is silently assigned the smooth formula's
$X^{-1}$ coefficient. The envelope construction is the rigorous discrete
inverse statement.

\section{Finite counts and exact reproducible checks}\label{sec:checks}
\subsection{An exact meet in the middle formula}
Let
\[
 V_n=\{v\in\{0,1\}^n:w^Tv=0\},\qquad b_n=|V_n|.
\]
Every admissible matrix is a sequence of rows in $V_n$. Split the row weights
into positive, negative, and, if $n$ is odd, one zero weight. Reflection pairs
the positive and negative lists. For a choice of rows $v_k\in V_n$ indexed by
the positive weights $k$, form $s=\sum_{k>0}kv_k\in\Z^n$ and let $f_n(s)$ be
the number of such choices producing $s$. Then
\begin{equation}\label{eq:MITM}
 C_n=b_n^{\,n\bmod2}\sum_{s\in\Z^n}f_n(s)^2.
\end{equation}
The column constraints say that the positive and reflected-negative weighted
row sums agree; their choices are independent and use the same list of
weights. When present, the row of weight zero is any element of $V_n$. This
proves \eqref{eq:MITM} exactly.

The finite values are
\begin{center}
\begin{tabular}{r r r r}
\toprule
$n$&$b_n$&$C_n$&Positive row tuples\\
\midrule
1&2&2&1\\
2&2&2&2\\
3&4&16&4\\
4&4&16&16\\
5&8&512&64\\
6&8&512&512\\
7&20&323600&8000\\
8&18&144930&104976\\
9&52&1808243216&7311616\\
\bottomrule
\end{tabular}
\end{center}
They agree with the displayed square array in OEIS A222959 \cite{oeis959}.
The independent finite enumeration underlying the last line uses $7{,}311{,}616$
positive weighted row tuples. The standard build's bounded check range and any
optional larger enumeration are documented explicitly in the code package;
recorded $n=9$ evidence is not mislabeled as a fresh default replay.

\subsection{Why these values do not validate the asymptotic numerically}
For $n=7,8,9$, the ratios of the exact counts to $2^{n^2}G_ne^{-7/5}$ are
approximately $21.15$, $3558.9$, and $23.66$. These large discrepancies are
strongly preasymptotic, not confirmation of numerical asymptotic accuracy.
A simple admissible sector already explains some small-dimension structure:
matrices palindromic in both rows and columns have
$\lceil n/2\rceil^2$ free binary entries, and hence
\begin{equation}\label{eq:palindromic}
 C_n\ge2^{\lceil n/2\rceil^2}.
\end{equation}
Opposite weights cancel pairwise in each row and column, proving this bound.
The sector is exponentially negligible relative to the asymptotic main count
as $n\to\infty$, since the difference in leading binary exponents is
$3n^2/4+O(n)$, dominating $O(n\log n)$.

Small row alphabets also have accidental rank deficiencies: through $n=14$,
their sizes are
\[
 2,2,4,4,8,8,20,18,52,48,152,138,472,428,
\]
and their real linear ranks are
\[
 1,1,2,2,3,3,5,5,8,9,10,11,12,13.
\]
In particular fixed-width patterns from the first few cases cannot replace a
growing-dimensional proof. These are exact finite observations, not hypotheses
used by the analytic argument.

\subsection{What the code establishes}
The accompanying source archive contains the manuscript, source ledger,
standard-library exact checks, recorded data with provenance, and a strict
reproducible build and verification workflow. The mathematical checks use
integer and rational arithmetic and explicit exceptions rather than Python
assertions, so optimized Python does not disable them. They check finite
projection identities, normalized moments, the first two coefficients, exact
structural identities, and bounded direct enumeration. A symbolic connected
Wick computation proves the two second-order contraction polynomials by exact
coefficient collection; a separately implemented raw Gaussian moment recursion
checks them for $n=2,3,4,5$. The builder and its
negative tests check the declared file inventory, manifest integrity, output
safety, and normal versus optimized replay. Exact ranges and commands are
recorded in the package documentation and machine-readable outputs.

Finite checking is corroboration of finite algebra and implementation. It does
not prove the global localization theorem, certify an eventual asymptotic
threshold, evaluate the uncomputed $c_j$ for $j\ge3$, or constitute formal
verification in a proof assistant. Reproducible PDF bytes are claimed only for
a fixed installed TeX toolchain; independent toolchain versions may render the
same source with different bytes.

\section{Primary source comparison and limits of the search}\label{sec:literature}
This section separates established predecessors from the specific theorem
proved here. The comparison is based on the primary statements identified
below, rather than titles alone. Its scope is a focused search completed on
4 October 2026; it is not exhaustive.

\subsection{The sequence and one dimensional fair words}
OEIS A222959 defines binary arrays with zero least-squares slope in each row
and column, and displays the finite square counts used above
\cite{oeis959}. The companion A222955 concerns the one-row condition
\cite{oeis955}. The displayed array was available, and its square entries
through nine were independently reproduced; the direct A222959 b-file was
not available in the source retrieval. That access limitation is not evidence
that the sequence lacks existing data.

Prodinger's 1979 paper studies binary words with equal counts of scattered
$ab$ and $ba$, relates them to a centered weighted sign sum, and proves
parity-sensitive one-dimensional asymptotics in Theorem~4, pp.~273--274
\cite{prodinger}. Richomme's later primary treatment gives the weighted-sum
characterization in Corollary~5.11, attributes the asymptotics to Prodinger in
Theorem~5.12, and proves the equivalence with zero least-squares slope in
Theorem~5.13 \cite{richomme}. In the notation of the row alphabet above, the
quoted one-dimensional leading forms are
\[
 b_{2m}\sim2^{2m-1}\sqrt{\frac3\pi}\,m^{-3/2},\qquad
 b_{2m+1}\sim2^{2m+1}\sqrt{\frac3\pi}\,m^{-3/2}.
\]
The row interpretation and its asymptotics are therefore prior work. They do
not by themselves count simultaneous row and column constraints in a square
matrix, whose number of independent equations grows with $n$.

\subsection{Fair picture arrays have weaker aggregate conditions}
Bera, Sriram, Nagar, Pan, and Subramanian study Parikh matrices of binary picture
arrays \cite{bera}. Their Definition~5 of fairness requires equal totals of
scattered $ab$ and $ba$ aggregated over all rows, and an analogous aggregate
condition over columns. The surrounding results, including Theorem~4 and
Sections~3--7, concern algebraic operations, concatenation, powers, and
geometric properties. Requiring each individual row and column to be fair is
stronger. For instance the $2\times2$ identity array has canceling aggregate
row and column imbalances, but neither of its length-two rows is individually
fair. This literature is relevant terminology and history, not a located
asymptotic enumeration theorem for the present count.

\subsection{Prescribed nullvectors and singular random matrices}
Arratia and DeSalvo develop lower-bound expansions for singularity of Bernoulli
sign matrices using nullvector templates associated with integer partitions
\cite{arratia}. Their definitions, Theorem~1, Proposition~9, Lemma~3, and the
full or almost-full support discussion in Section~5 address prescribed
nullvector events and their interactions. The fixed-template constants in
Proposition~9 depend on the template support and one-row concentration.
Propositions~10--11 condition a prescribed nullvector against exclusions such
as sparse opposite-side templates; they do not count two prescribed growing
arithmetic-progression nullvectors simultaneously. The present inference is
limited: those estimates cannot be silently made uniform when both the
support and the coefficients grow with $n$. They are conceptual predecessors,
not an inspected matching count equivalent.

\subsection{Maximum entropy Gaussian and Edgeworth frameworks}
The binary counting theorem, Theorem~2.6, of Barvinok and Hartigan's Gaussian
paper has quantitative covariance and boundedness hypotheses
\cite{bhgaussian}. Its parameter $\lambda$ is a lower bound for the quadratic Fourier form
$q(t)=\tfrac12t^T\operatorname{Cov}\,t$. At fixed relative accuracy it must
dominate a constant times $\theta^2d^2\log N$, where $\theta$ controls the
constraint-column $\ell^1$ norms. The natural saturated coordinates here fail
this sufficient condition. Delete one column-margin coordinate whose weight
is $1$; then $d=2n-1$, $N=n^2$, and $\theta=\Omega(n)$. Remove that coordinate
from the full gauge vector $(w,-w)$. The resulting vector has squared norm
$2S-1$, and its image under the reduced design transpose has squared norm $S$.
The Bernoulli covariance is one fourth of the reduced design Gram matrix, so
\[
 \lambda_{\min}(\operatorname{Cov})\le\frac{S}{4(2S-1)},\qquad
 \lambda_{\min}(q)\le\frac{S}{8(2S-1)}=O(1).
\]
Thus the required domination fails in these coordinates. The factor of two
between the bounds is solely the source's quadratic-form normalization. No
claim is made about every possible reformulation. The paper also explains
that these sufficient conditions do not establish ordinary two-way-table
Gaussian asymptotics.

The later Edgeworth paper states a general program with explicit conditions
I--V in Theorem~1, pp.~5--7 \cite{bhedgeworth}. Condition~V is precisely a
relative negligibility requirement for the Fourier integral outside a small
central region. Its ordinary contingency-margin specialization does not
verify that condition for growing signed arithmetic-progression coefficients.
We credit the maximum-entropy, Fourier, and Wick correction framework; we
prove the necessary problem-specific global estimate in
Sections~\ref{sec:inverselemmas}--\ref{sec:defects} instead of treating the
framework as an automatic theorem.

\subsection{General high dimensional lattice counting}
Kuperberg, Lovett, and Peled give a general local-limit and counting framework
for regular combinatorial structures \cite{klp}, including Theorems~2.4--2.5
and the basis-free Theorem~4.15. One hypothesis requires a coordinate
permutation symmetry group acting transitively on the ground set. Here the
ground set consists of the $n^2$ cells, and the constraint space is $E$.
A coordinate permutation preserving $E$ also preserves its orthogonal
projection diagonal. But $h_{ij}=a_i+a_j-a_ia_j$ varies across cells for large
$n$, ruling out that transitivity. Their counting theorem's size lower bound
of order $\dim(E)^6$, even before boundedness and logarithmic factors, also
exceeds $n^2$ in this regime. These actual hypothesis obstructions explain why
that important general precedent does not directly establish our theorem.

\subsection{Other nearby matrix models}
Harrison and Miller's weighted binary matrix model assigns nonuniform cell odds
to matrices with ordinary prescribed row and column sums \cite{harrison}.
Its weights are probabilities, not growing signed coefficients inside each
margin equation. Definition~(1) and Section~4.2 establish the different scope:
importance sampling with ordinary margins and associated approximations.

Aw's admixed arrays involve paired binary matrices with uniform unweighted
row and paired-column sums \cite{aw}. The inspected model, Theorem~4.3 setup,
reduced cosine-squared integral, and local expansion use coefficients one in
the phase constraints. The local quartic approach is relevant, but no general
growing-weight global-localization theorem is imported from that setting.
In particular, a modular cosine bound is not extended to unwrapped real
phases outside a central interval.

\subsection{Scoped conclusion}
The focused comparison did not locate a primary theorem matching the
simultaneous centered-arithmetic-progression left and right nullvector count,
the factor $e^{-7/5}$, or its rational all-fixed-orders expansion. This is a
limited negative search result. The safe mathematical contribution stated
here is a complete application and extension of established Fourier and
higher-order Gaussian enumeration methods, including a fully explicit
defect-core localization proof. Absolute priority remains open. The package's
source ledger records links, inspected theorem scope, and retrieval limits;
no novelty assertion rests on a search-result title or absence from a catalog.

\section{Limitations and further research}\label{sec:limits}
The following questions are not resolved by the theorem or by its code.
\begin{enumerate}
\item \textbf{Practical error and onset.} The proof is eventual and deliberately
uses tiny absolute thresholds. A useful explicit error bound at accessible
$n$ would require substantially sharper localization and local estimates.
The small exact values cannot locate that onset.
\item \textbf{Higher explicit coefficients.} Formula~\eqref{eq:cj} gives every
fixed coefficient by a finite computation, and $c_1,c_2$ have been evaluated
here. The supplied contraction-class computation treats the two contractions
needed at second order. Extending that implementation systematically could
make $c_3,c_4,\ldots$ practical; those coefficients are not supplied as
computed results.
\item \textbf{Optimal truncation.} Bounds uniform in a growing order $J$,
convergence, Borel summability, and exponentially improved expansions all lie
beyond this fixed-order argument.
\item \textbf{Other weight lists and rectangular limits.} The lattice and
Gaussian core formulas are rectangular, but the present global theorem uses
both dense arithmetic-progression circuits and square-scale defect costs.
Extending it to rectangular aspect ratios or other deterministic nullvectors
requires checking new arithmetic stability and localization hypotheses.
\item \textbf{Transition from exceptional finite structure.} The exact
palindromic sector and row-alphabet rank defects suggest a useful study of
intermediate dimensions. They explain caution, but do not supply a rigorous
crossover law.
\item \textbf{Integer inverse boundaries.} The envelopes shrink to one odd
integer except near threshold coincidences. Sharper asymptotics do not alone
remove every possible rounding ambiguity; exact comparison or explicit
remainder bounds are needed there.
\item \textbf{Literature and formalization.} A broader primary-source search
could find additional predecessors. A formal proof assistant development
would require encoding the global inverse lemmas, Haar normalization, and
uniform analytic remainders; the current standard-library checks are not
such a formalization.
\end{enumerate}

\clearpage
\begin{thebibliography}{99}
\bibitem{oeis959} OEIS Foundation Inc., \emph{A222959}, binary arrays with every
row and column having zero least-squares slope,
\url{https://oeis.org/A222959}. Accessed 4 October 2026.
\bibitem{oeis955} OEIS Foundation Inc., \emph{A222955}, the associated binary
word count, \url{https://oeis.org/A222955}. Accessed 4 October 2026.
\bibitem{prodinger} H. Prodinger, \emph{On a generalization of the Dyck-language
over a two letter alphabet}, Discrete Mathematics \textbf{28} (1979), 269--276.
\url{https://doi.org/10.1016/0012-365X(79)90134-1}.
Author copy: \url{https://www.finanz.math.tugraz.at/~prodinger/dyck.pdf}.
\bibitem{richomme} G. Richomme, \emph{On some 2-binomial coefficients of binary
words: geometrical interpretation, partitions of integers, and fair words},
arXiv:2510.07159v1 (2025), \url{https://arxiv.org/html/2510.07159v1}.
\bibitem{bera} S. Bera, S. Sriram, A. K. Nagar, L. Pan, and K. G. Subramanian,
\emph{Algebraic Properties of Parikh Matrices of Binary Picture Arrays},
Journal of Mathematics (2020), article 3236405,
\url{https://doi.org/10.1155/2020/3236405}.
\bibitem{arratia} R. Arratia and S. DeSalvo, \emph{On the singularity of random
Bernoulli matrices---novel integer partitions and lower bound expansions},
arXiv:1105.2834v2 (2012), \url{https://arxiv.org/pdf/1105.2834}.
\bibitem{bhgaussian} A. Barvinok and J. A. Hartigan, \emph{Maximum entropy
Gaussian approximation for the number of integer points and volumes of
polytopes}, arXiv:0903.5223v2 (2009), \url{https://arxiv.org/pdf/0903.5223}.
\bibitem{bhedgeworth} A. Barvinok and J. A. Hartigan, \emph{Maximum entropy
Edgeworth estimates of the number of integer points in polytopes},
arXiv:0910.2497v2 (2010), \url{https://arxiv.org/pdf/0910.2497}.
\bibitem{klp} G. Kuperberg, S. Lovett, and R. Peled, \emph{Probabilistic existence
of regular combinatorial structures}, arXiv:1302.4295v3 (2017),
\url{https://arxiv.org/pdf/1302.4295}.
\bibitem{harrison} M. T. Harrison and J. W. Miller, \emph{Importance sampling
for weighted binary random matrices with specified margins},
arXiv:1301.3928v1 (2013), author copy
\url{https://jwmi.github.io/publications/Nonuniform_matrices.pdf}.
\bibitem{aw} A. J. Aw, \emph{Asymptotic enumeration of admixed arrays and a
different independence heuristic}, arXiv:2604.08857v1 (2026),
\url{https://arxiv.org/html/2604.08857v1}.
\end{thebibliography}
\end{document}
